\documentclass{amsart}
% For dealing with references we use the comment environment
\usepackage{verbatim}
\newenvironment{reference}{\comment}{\endcomment}
%\newenvironment{reference}{}{}
\newenvironment{slogan}{\comment}{\endcomment}
\newenvironment{history}{\comment}{\endcomment}

% For commutative diagrams we use Xy-pic
\usepackage[all]{xy}

% We use 2cell for 2-commutative diagrams.
\xyoption{2cell}
\UseAllTwocells

% We use multicol for the list of chapters between chapters
\usepackage{multicol}

% This is generally recommended for better output
\usepackage{lmodern}
\usepackage[T1]{fontenc}

% For cross-file-references

% Package for hypertext links:
\usepackage{hyperref}

% For any local file, say "hello.tex" you want to link to please

% Theorem environments.
%
\theoremstyle{plain}
\newtheorem{theorem}[subsection]{Theorem}
\newtheorem{proposition}[subsection]{Proposition}
\newtheorem{lemma}[subsection]{Lemma}

\theoremstyle{definition}
\newtheorem{definition}[subsection]{Definition}
\newtheorem{example}[subsection]{Example}
\newtheorem{exercise}[subsection]{Exercise}
\newtheorem{situation}[subsection]{Situation}

\theoremstyle{remark}
\newtheorem{remark}[subsection]{Remark}
\newtheorem{remarks}[subsection]{Remarks}

\numberwithin{equation}{subsection}

% Macros
%
\def\lim{\mathop{\mathrm{lim}}\nolimits}
\def\colim{\mathop{\mathrm{colim}}\nolimits}
\def\Spec{\mathop{\mathrm{Spec}}}
\def\Hom{\mathop{\mathrm{Hom}}\nolimits}
\def\Ext{\mathop{\mathrm{Ext}}\nolimits}
\def\SheafHom{\mathop{\mathcal{H}\!\mathit{om}}\nolimits}
\def\SheafExt{\mathop{\mathcal{E}\!\mathit{xt}}\nolimits}
\def\Sch{\mathit{Sch}}
\def\Mor{\mathop{\mathrm{Mor}}\nolimits}
\def\Ob{\mathop{\mathrm{Ob}}\nolimits}
\def\Sh{\mathop{\mathit{Sh}}\nolimits}
\def\NL{\mathop{N\!L}\nolimits}
\def\CH{\mathop{\mathrm{CH}}\nolimits}
\def\proetale{{pro\text{-}\acute{e}tale}}
\def\etale{{\acute{e}tale}}
\def\QCoh{\mathit{QCoh}}
\def\Ker{\mathop{\mathrm{Ker}}}
\def\Im{\mathop{\mathrm{Im}}}
\def\Coker{\mathop{\mathrm{Coker}}}
\def\Coim{\mathop{\mathrm{Coim}}}

% Boxtimes
%
\DeclareMathSymbol{\boxtimes}{\mathbin}{AMSa}{"02}

%
% Macros for moduli stacks/spaces
%
\def\QCohstack{\mathcal{QC}\!\mathit{oh}}
\def\Cohstack{\mathcal{C}\!\mathit{oh}}
\def\Spacesstack{\mathcal{S}\!\mathit{paces}}
\def\Quotfunctor{\mathrm{Quot}}
\def\Hilbfunctor{\mathrm{Hilb}}
\def\Curvesstack{\mathcal{C}\!\mathit{urves}}
\def\Polarizedstack{\mathcal{P}\!\mathit{olarized}}
\def\Complexesstack{\mathcal{C}\!\mathit{omplexes}}
% \Pic is the operator that assigns to X its picard group, usage \Pic(X)
% \Picardstack_{X/B} denotes the Picard stack of X over B
% \Picardfunctor_{X/B} denotes the Picard functor of X over B
\def\Pic{\mathop{\mathrm{Pic}}\nolimits}
\def\Picardstack{\mathcal{P}\!\mathit{ic}}
\def\Picardfunctor{\mathrm{Pic}}
\def\Deformationcategory{\mathcal{D}\!\mathit{ef}}

\setlength{\emergencystretch}{3em}
\begin{document}
\title[Stacks Project, Tag 02KK]{628 --- Finite presentation descends through a faithfully flat finitely presented algebra}
\maketitle
\noindent Copyright (C) 2005--2025 Johan de Jong. Stacks Project authors. Source: \href{https://stacks.math.columbia.edu/tag/02KK}{Tag 02KK}.

\noindent Editorial difficulty note (not author text). Difficulty H2: The argument spreads a faithfully flat algebra to a filtered stage, controls eventual surjectivity through its open image, and descends a compatible equalizer diagram. Managing finite-presentation approximation and faithful flatness simultaneously is beyond standard qualification exercises..

\noindent Editorial provenance note (not author text). History: excerpt from revision \texttt{a0444\allowbreak 6e57e\allowbreak c1fbc\allowbreak 252a8\allowbreak 71afc\allowbreak ec775\allowbreak 2fb28\allowbreak 07b14}, descent.tex, lines 4260--4345. Reference presentation was adapted; original-author context and core support blocks were retained or added. The mathematical wording is unchanged. Licensed under GNU FDL 1.2 or later; no invariant sections or cover texts. See ../licenses/COPYING. Context blocks reproduce author definitions and supporting results; they are not additional counted problems.

\begin{lemma}
\label{lemma-flat-finitely-presented-permanence-algebra}
Let $R \to A \to B$ be ring maps.
Assume $R \to B$ is of finite presentation and
$A \to B$ faithfully flat and of finite presentation.
Then $R \to A$ is of finite presentation.
\end{lemma}

\begin{proof}
Consider the algebra $C = B \otimes_A B$ together with the
pair of maps $p, q : B \to C$ given by $p(b) = b \otimes 1$
and $q(b) = 1 \otimes b$. Of course the two compositions
$A \to B \to C$ are the same. Note that as
$p : B \to C$ is flat and of finite presentation (base change of
$A \to B$), the ring map $R \to C$ is of finite presentation
(as the composite of $R \to B \to C$).

\medskip\noindent
We are going to use the criterion
Algebra, Lemma \href{https://stacks.math.columbia.edu/tag/00QO}{lemma characterize finite presentation (Tag 00QO)}
to show that $R \to A$ is of finite presentation.
Let $S$ be any $R$-algebra, and suppose that
$S = \colim_{\lambda \in \Lambda} S_\lambda$ is written
as a directed colimit of $R$-algebras.
Let $A \to S$ be an $R$-algebra homomorphism. We have to
show that $A \to S$ factors through one of the $S_\lambda$.
Consider the rings $B' = S \otimes_A B$ and
$C' = S \otimes_A C = B' \otimes_S B'$.
As $B$ is faithfully flat of finite presentation over $A$, also $B'$
is faithfully flat of finite presentation over $S$.
By Algebra, Lemma \href{https://stacks.math.columbia.edu/tag/02JO}{lemma flat finite presentation limit flat (Tag 02JO)}
part (2) applied to the pair $(S \to B', B')$ and the system $(S_\lambda)$
there exists a $\lambda_0 \in \Lambda$
and a flat, finitely presented $S_{\lambda_0}$-algebra
$B_{\lambda_0}$ such that $B' = S \otimes_{S_{\lambda_0}} B_{\lambda_0}$.
For $\lambda \geq \lambda_0$ set
$B_\lambda = S_\lambda \otimes_{S_{\lambda_0}} B_{\lambda_0}$ and
$C_\lambda = B_\lambda \otimes_{S_\lambda} B_\lambda$.

\medskip\noindent
We interrupt the flow of the argument to show that $S_\lambda \to B_\lambda$
is faithfully flat for $\lambda$ large enough. (This should really
be a separate lemma somewhere else, maybe in the chapter on limits.)
Since $\Spec(B_{\lambda_0}) \to \Spec(S_{\lambda_0})$ is
flat and of finite presentation it is open (see Morphisms,
Lemma \href{https://stacks.math.columbia.edu/tag/01UA}{lemma fppf open (Tag 01UA)}).
Let $I \subset S_{\lambda_0}$ be an ideal such that
$V(I) \subset \Spec(S_{\lambda_0})$ is the complement
of the image. Note that formation of the image commutes
with base change. Hence, since $\Spec(B') \to \Spec(S)$
is surjective, and $B' = B_{\lambda_0} \otimes_{S_{\lambda_0}} S$
we see that $IS = S$. Thus for some $\lambda \geq \lambda_0$ we
have $IS_{\lambda} = S_\lambda$. For this and all greater
$\lambda$ the morphism
$\Spec(B_\lambda) \to \Spec(S_\lambda)$ is surjective.

\medskip\noindent
By analogy with the notation in the first paragraph of the proof denote
$p_\lambda, q_\lambda : B_\lambda \to C_\lambda$ the two canonical maps.
Then $B' = \colim_{\lambda \geq \lambda_0} B_\lambda$
and $C' = \colim_{\lambda \geq \lambda_0} C_\lambda$.
Since $B$ and $C$ are finitely presented over $R$ there exist
(by Algebra, Lemma \href{https://stacks.math.columbia.edu/tag/00QO}{lemma characterize finite presentation (Tag 00QO)}
applied several times)
a $\lambda \geq \lambda_0$ and an $R$-algebra maps
$B \to B_\lambda$, $C \to C_\lambda$ such that
the diagram
$$
\xymatrix{
C \ar[rr] & &
C_\lambda \\
B \ar[rr]
\ar@<1ex>[u]^-p
\ar@<-1ex>[u]_-q
& &
B_\lambda
\ar@<1ex>[u]^-{p_\lambda}
\ar@<-1ex>[u]_-{q_\lambda}
}
$$
is commutative. OK, and this means that $A \to B \to B_\lambda$
maps into the equalizer of $p_\lambda$ and $q_\lambda$.
By Lemma \href{https://stacks.math.columbia.edu/tag/023M}{lemma ff exact (Tag 023M)} we
see that $S_\lambda$ is the equalizer of $p_\lambda$ and $q_\lambda$.
Thus we get the desired ring map $A \to S_\lambda$ and we win.
\end{proof}
\end{document}
